% Ch2 WS1 -- atomic theory, laws, isotope notation. Blocks 1-2.
\documentclass{apchem-worksheet}

\apcunitword{Chapter}
\apcunit{2}
\apcunittitle{Atoms, Molecules, and Ions}
\apcdoctitle{Worksheet 1 \textbullet\ Laws, Experiments, Isotopes}
\apcsubtitle{Zumdahl \S2.1--2.7 \textbullet\ evidence-to-conclusion reasoning}

\begin{document}
\wsheader[For every historical item, name the evidence before the
conclusion. ``Rutherford found the nucleus'' earns nothing without the
observation that forced it.]

\problem[]{Two compounds of sulfur and oxygen are analyzed. In compound A,
\SI{1.00}{\gram} of S combines with \SI{1.00}{\gram} of O. In compound B,
\SI{1.00}{\gram} of S combines with \SI{1.50}{\gram} of O.}
\begin{parts}
  \item Compute the ratio of oxygen masses per gram of sulfur.
        \answerblank[3cm]{$1.50/1.00 = 3{:}2$}
  \item Which law does this illustrate?
        \answerblank[5.6cm]{law of multiple proportions}
  \item Given A is \ce{SO2}, propose a formula for B and justify.
        \answerlines[2]{\ce{SO3} --- if A has 2 O per S, a 3:2 oxygen ratio
        gives B 3 O per S.}
\end{parts}

\problem[]{Match each observation to the law it demonstrates:}
\begin{parts}
  \item A sealed flask weighs the same before and after a reaction inside
        it. \answerblank[5.5cm]{conservation of mass}
  \item Carbon dioxide from a volcano and from combustion both contain
        27.3\% C. \answerblank[5.5cm]{definite proportion}
\end{parts}

\problem[]{Thomson's cathode-ray tube.}
\begin{parts}
  \item Which way does the beam deflect in an electric field, and what does
        that establish? \answerlines[2]{Toward the positive plate --- the
        particles carry negative charge.}
  \item Thomson measured the same charge-to-mass ratio using several
        different gases. State the conclusion this supports.
        \answerlines[2]{The same negatively charged particle is present in
        all elements --- electrons are universal components of atoms.}
  \item Thomson's model placed electrons in a diffuse positive sphere. Which
        later experiment refuted it, and by what specific observation?
        \answerlines[2]{Rutherford's gold foil --- a small fraction of
        $\alpha$ particles deflected at large angles or straight back, which
        a diffuse positive charge could never cause.}
\end{parts}

\problem[]{Complete the table.}

\renewcommand{\arraystretch}{1.6}
\noindent\begin{tabular}{@{}lcccc@{}}
\toprule
\textbf{Species} & \textbf{protons} & \textbf{neutrons} & \textbf{electrons} & \textbf{mass \#}\\
\midrule
\ce{^{31}_{15}P} & \answerblank[1.1cm]{15} & \answerblank[1.1cm]{16} &
  \answerblank[1.1cm]{15} & \answerblank[1.1cm]{31}\\
\ce{^{81}_{35}Br^-} & \answerblank[1.1cm]{35} & \answerblank[1.1cm]{46} &
  \answerblank[1.1cm]{36} & \answerblank[1.1cm]{81}\\
\ce{^{52}_{24}Cr^{3+}} & \answerblank[1.1cm]{24} & \answerblank[1.1cm]{28} &
  \answerblank[1.1cm]{21} & \answerblank[1.1cm]{52}\\
\answerblank[1.8cm]{\ce{^{16}O^{2-}}} & 8 & 8 & 10 & 16\\
\bottomrule
\end{tabular}

\problem[]{An ion has 38 protons, 50 neutrons, and 36 electrons.}
\begin{parts}
  \item Write its complete isotope symbol.
        \answerblank[3.5cm]{\ce{^{88}_{38}Sr^{2+}}}
  \item Name one species isoelectronic with it.
        \answerblank[4.6cm]{\ce{Kr}, \ce{Br-}, \ce{Rb+}, \ce{Y^3+}\ldots}
\end{parts}

\problem[]{Predict the ion each element forms, then the formula of the
compound it makes with the partner given:}
\begin{parts}
  \item Ca with Br: \answerblank[5cm]{\ce{Ca^2+}, \ce{Br-} $\to$ \ce{CaBr2}}
  \item Li with N: \answerblank[5cm]{\ce{Li+}, \ce{N^3-} $\to$ \ce{Li3N}}
  \item Al with S: \answerblank[5cm]{\ce{Al^3+}, \ce{S^2-} $\to$ \ce{Al2S3}}
\end{parts}

\begin{spiralstrip}{prerequisites}
  \item Sig figs in \SI{0.04500}{\gram}: \answerblank[2cm]{4}
  \item \SI{2.5}{\kilo\gram} in grams: \answerblank[2.5cm]{\SI{2500}{\gram}}
  \item Which is a mixture: air, \ce{CO2}, or copper?
        \answerblank[2.5cm]{air}
  \item Protons in an atom with $Z = 47$: \answerblank[2cm]{47}
  \item Isotopes differ in what particle? \answerblank[2.5cm]{neutrons}
\end{spiralstrip}

\end{document}
